% =====================================================================
%  preamble.tex -- shared preamble for the standard manuscript style.
%  Derived from the dominant conventions of the openai/math preprint set:
%  \documentclass[11pt]{article}, modular \input of sections/, one shared
%  theorem counter, BibTeX with \bibliographystyle{plain}.
% =====================================================================

% ---------- 1. Encoding, fonts, page geometry ----------
\usepackage[T1]{fontenc}
\usepackage[utf8]{inputenc}
\usepackage{lmodern}
\usepackage[margin=1in]{geometry}

% ---------- 2. Mathematics ----------
\usepackage{amsmath,amssymb,amsthm,mathtools}

% ---------- 3. Layout, lists, tables, figures ----------
\usepackage{needspace}
\usepackage{enumitem,booktabs,array}
\usepackage{microtype}
\usepackage{tikz}
\usetikzlibrary{arrows.meta,positioning,calc,decorations.pathreplacing}
\usepackage{graphicx}
\usepackage{xcolor}

% ---------- 4. Cross-references and links ----------
\definecolor{linkblue}{RGB}{25,65,105}
\usepackage[colorlinks=true,linkcolor=linkblue,citecolor=linkblue,
            urlcolor=linkblue,breaklinks=true]{hyperref}
\usepackage[capitalise,noabbrev,nameinlink]{cleveref}

% ---------- 5. Numbering: equations and theorems ----------
\numberwithin{equation}{section}

% "theorem" owns the counter and is numbered within the section;
% every other environment shares it, so results read Theorem 2.1,
% Lemma 2.2, Proposition 2.3, ... in one running sequence.
\newtheorem{theorem}{Theorem}[section]
\newtheorem{lemma}[theorem]{Lemma}
\newtheorem{proposition}[theorem]{Proposition}
\newtheorem{corollary}[theorem]{Corollary}
\theoremstyle{definition}
\newtheorem{definition}[theorem]{Definition}
\newtheorem{example}[theorem]{Example}
\newtheorem{convention}[theorem]{Convention}
\newtheorem{assumption}[theorem]{Assumption}
\newtheorem{construction}[theorem]{Construction}
\theoremstyle{remark}
\newtheorem{remark}[theorem]{Remark}
\newtheorem{notation}[theorem]{Notation}
\newtheorem{claim}[theorem]{Claim}

% ---------- 6. Page-break hygiene ----------
% Keep a few lines together before subsections and before each
% theorem-like environment, so a statement never starts on the last
% line of a page.
\AddToHook{cmd/subsection/before}{\Needspace{9\baselineskip}}
\AddToHook{env/theorem/before}{\Needspace{2\baselineskip}}
\AddToHook{env/lemma/before}{\Needspace{2\baselineskip}}
\AddToHook{env/proposition/before}{\Needspace{2\baselineskip}}
\AddToHook{env/corollary/before}{\Needspace{2\baselineskip}}
\AddToHook{env/definition/before}{\Needspace{2\baselineskip}}
\AddToHook{env/remark/before}{\Needspace{2\baselineskip}}
\allowdisplaybreaks[2]

% ---------- 7. Lists and table of contents ----------
\setlist{itemsep=2pt,topsep=5pt}
\setcounter{tocdepth}{2}

% ---------- 8. Project-specific macros ----------
% Number sets. Extend or prune freely.
\newcommand{\N}{\mathbb{N}}
\newcommand{\Z}{\mathbb{Z}}
\newcommand{\Q}{\mathbb{Q}}
\newcommand{\R}{\mathbb{R}}
\newcommand{\C}{\mathbb{C}}
\newcommand{\F}{\mathbb{F}}

% Operators.
\DeclareMathOperator{\rank}{rank}
\DeclareMathOperator{\supp}{supp}
\DeclareMathOperator{\im}{im}
\DeclareMathOperator{\id}{id}
\DeclareMathOperator{\tr}{tr}
\DeclareMathOperator{\Spec}{Spec}
\DeclareMathOperator{\Hom}{Hom}
\DeclareMathOperator{\End}{End}

% Utility macros used across the sample sections.
\newcommand{\dd}{\mathrm{d}}
\newcommand{\abs}[1]{\left\lvert #1 \right\rvert}
\newcommand{\norm}[1]{\left\lVert #1 \right\rVert}
\newcommand{\ip}[2]{\left\langle #1, #2 \right\rangle}
\newcommand{\emphterm}[1]{\emph{#1}}
